% !TEX program = xelatex
% 2026 年高教社杯全国大学生数学建模竞赛 A 题
% 问题二、三、四：整个烘干过程的耦合传热传质模型
%
% 独立可编译文档；并入整篇论文时取 \section 及其后全部内容即可。

\documentclass[12pt,a4paper]{ctexart}

\usepackage{amsmath,amssymb,bm}
\usepackage{booktabs}
\usepackage{multirow}
\usepackage{array}
\usepackage{geometry}
\usepackage{caption}

\geometry{left=2.7cm,right=2.7cm,top=2.6cm,bottom=2.6cm}
\captionsetup{font=small,labelsep=quad}
\linespread{1.32}

\begin{document}

\section{整个烘干过程的耦合传热传质模型（问题二）}

\subsection{问题分析}

问题二要求建立\emph{整个}烘干过程（预热平衡 $+$ 恒温干燥，持续 $2\sim3$\,天）
药材温度与水分浓度的变化规律。与问题一的本质区别在于物性不再恒定：
附录 3 给出
\begin{equation}
  \rho=650+128C,\qquad
  c_p=1450+2736\frac{C}{C+1},\qquad
  k=0.21+0.38\frac{C}{C+1},
  \label{eq:p2-props}
\end{equation}
\begin{equation}
  D=2.4\times10^{-3}\,e^{-0.45/C}e^{-3850/T},
  \label{eq:p2-D}
\end{equation}
其中 $D$ 同时依赖含水率 $C$ 与温度 $T$（$T$ 取绝对温标 K），
$\rho$ 也随 $C$ 变化。因此水分方程中的扩散系数随温度演化，
而温度方程中的热容、导热系数随含水率演化，\textbf{两个方程必须联立求解}，
问题一中"热--质解耦"的结构性质不再成立。

\subsection{模型假设}

在问题一假设的基础上补充与修改：

\begin{enumerate}
  \item[(B1)] 预热平衡阶段与恒温干燥阶段使用同一组物性关系式
        \eqref{eq:p2-props}、\eqref{eq:p2-D}，两阶段的差别仅体现在
        烘房工况 $T_\infty$、$C_\infty$ 上；
  \item[(B2)] 恒温干燥阶段烘房工况恒定，取 $T_\infty=50\,^\circ$C、
        $C_\infty=0.05$\,kg/kg（依据见 §\ref{sec:bc}）；
  \item[(B3)] 问题二、三中不考虑药材尺寸变化（收缩由问题四处理），
        半径恒为 $R=2$\,cm；
  \item[(B4)] 对流换热系数与对流传质系数题面未给出，
        沿用附录 2 的 $h=25$\,W/(m$^2\cdot$K)、$h_m=8\times10^{-7}$\,m/s；
  \item[(B5)] 忽略汽化潜热，理由同问题一。
\end{enumerate}

\subsection{模型建立}

\subsubsection{控制方程}

圆柱坐标下的一维径向守恒方程：
\begin{equation}
  \frac{\partial(\rho c_p T)}{\partial t}
  =\frac{1}{r}\frac{\partial}{\partial r}
   \left(k\,r\frac{\partial T}{\partial r}\right),
  \label{eq:p2-heat}
\end{equation}
\begin{equation}
  \frac{\partial(\rho C)}{\partial t}
  =\frac{1}{r}\frac{\partial}{\partial r}
   \left(\rho D\,r\frac{\partial C}{\partial r}\right).
  \label{eq:p2-mass}
\end{equation}

关于式 \eqref{eq:p2-mass} 的写法需要说明：附录 2 中 $\rho$ 是与 $C$ 无关的常数，
故问题一中 $\rho$ 在等式两侧约去，水分方程退化为简单扩散形式；
而附录 3、4 把 $\rho$ 写成 $C$ 的函数，说明它确实参与方程，
因此这里保留完整的守恒形式。

\subsubsection{定解条件}

\begin{equation}
  T(r,0)=28\,^\circ\mathrm{C},\qquad C(r,0)=2.55\ \mathrm{kg/kg},
  \label{eq:p2-ic}
\end{equation}
\begin{equation}
  \left.\frac{\partial T}{\partial r}\right|_{r=0}=0,\qquad
  \left.\frac{\partial C}{\partial r}\right|_{r=0}=0,
  \label{eq:p2-center}
\end{equation}
\begin{equation}
  -k\left.\frac{\partial T}{\partial r}\right|_{r=R}
  =h\left(T_s-T_\infty(t)\right),
  \label{eq:p2-bcT}
\end{equation}
\begin{equation}
  -D\left.\frac{\partial C}{\partial r}\right|_{r=R}
  =h_m\left(C_s-C_\infty(t)\right).
  \label{eq:p2-bcC}
\end{equation}

\subsubsection{烘房工况的确定}
\label{sec:bc}

附件 1 只覆盖 $0\sim14400$\,s。对其末段 1\,h（$10800\sim14400$\,s）取平均，
得 $T_\infty=49.9989\,^\circ$C、$C_\infty=0.04999$\,kg/kg —— 与整数
$50\,^\circ$C、$0.05$\,kg/kg 几乎完全一致，说明这正是题目设定的恒温干燥工况。
因此边界条件取分段形式：
\begin{equation}
  T_\infty(t)=
  \begin{cases}
    \text{附件 1 线性插值}, & 0\le t\le 14400\ \mathrm{s},\\[2pt]
    50\,^\circ\mathrm{C}, & t>14400\ \mathrm{s},
  \end{cases}
  \label{eq:p2-Tinf}
\end{equation}
$C_\infty(t)$ 同理，$t>14400$\,s 后取 $0.05$\,kg/kg。

\subsection{模型求解}

\subsubsection{空间离散}

仍采用节点中心有限体积法，但此时每个控制体的"热容"与"湿容"都随解变化。
对水分方程，左端按链式法则展开
\begin{equation}
  \frac{\partial(\rho C)}{\partial t}
  =\left(\rho+C\frac{\mathrm{d}\rho}{\mathrm{d}C}\right)\frac{\partial C}{\partial t},
  \qquad
  \rho+C\frac{\mathrm{d}\rho}{\mathrm{d}C}=650+256C .
  \label{eq:p2-accum}
\end{equation}
该系数与扩散通量系数 $\rho D$ 一并按上一迭代步的值滞后处理。

\subsubsection{时间离散与联立迭代}

时间方向用后向 Euler，非线性用 Picard 迭代：
\begin{equation}
  \begin{cases}
    \left(\bm{I}+\bm{A}_C^{(k)}\right)\bm{C}^{(k+1)}
      =\bm{b}\big(\bm{C}^{n},\bm{T}^{(k)}\big),\\[3pt]
    \left(\bm{I}+\bm{A}_T^{(k)}\right)\bm{T}^{(k+1)}
      =\bm{b}\big(\bm{T}^{n},\bm{C}^{(k+1)}\big),
  \end{cases}
  \qquad
  \big\|\cdot^{(k+1)}-\cdot^{(k)}\big\|_\infty<\varepsilon .
  \label{eq:p2-picard}
\end{equation}
取 $\varepsilon=10^{-11}$，通常 $2\sim3$ 次迭代收敛。每一步中两个方程各化为
一个三对角方程组，用追赶法在 $O(N)$ 时间内求解。

需要强调的一点：式 \eqref{eq:p2-picard} 右端的时间层 $n$ 必须是\textbf{本时间步
开始时}的状态，且在 Picard 迭代中保持不变；若误用上一迭代步的结果，
则每迭代一次就相当于再推进一个 $\Delta t$，数值解会比真实解快数倍。

\subsubsection{计算参数}

取 $N=200$（$\xi$ 方向）、$\Delta t=5$\,s，求解至中心含水率低于 $0.15$\,kg/kg。
收敛性检验见 §\ref{sec:p23-verify}。

\subsection{结果与分析}

表 \ref{tab:p2-T} 与表 \ref{tab:p2-C} 分别给出 3\,h 内每 $0.5$\,h 的温度与
水分浓度。

\begin{table}[htbp]
  \centering
  \caption{3 小时内药材的温度（单位：$^\circ$C）}
  \label{tab:p2-T}
  \begin{tabular}{cccccc}
    \toprule
    \multirow{2}{*}{时间/h} & \multicolumn{5}{c}{到药材中心的距离/cm} \\
    \cmidrule(lr){2-6}
     & 0     & 0.5   & 1.0   & 1.5   & 2.0 \\
    \midrule
    0.5 & 32.1881 & 32.3802 & 32.9617 & 33.9525 & 35.4034 \\
    1.0 & 40.3372 & 40.5087 & 41.0149 & 41.8338 & 42.9593 \\
    1.5 & 45.7798 & 45.8690 & 46.1307 & 46.5486 & 47.0929 \\
    2.0 & 48.3951 & 48.4342 & 48.5483 & 48.7294 & 48.9672 \\
    2.5 & 49.4342 & 49.4473 & 49.4842 & 49.5406 & 49.6414 \\
    3.0 & 49.8336 & 49.8397 & 49.8599 & 49.8970 & 49.9555 \\
    \bottomrule
  \end{tabular}
\end{table}

\begin{table}[htbp]
  \centering
  \caption{3 小时内药材的水分浓度（单位：kg/kg）}
  \label{tab:p2-C}
  \begin{tabular}{cccccc}
    \toprule
    \multirow{2}{*}{时间/h} & \multicolumn{5}{c}{到药材中心的距离/cm} \\
    \cmidrule(lr){2-6}
     & 0     & 0.5   & 1.0   & 1.5   & 2.0 \\
    \midrule
    0.5 & 2.5500 & 2.5498 & 2.5398 & 2.3807 & 1.6233 \\
    1.0 & 2.5435 & 2.5272 & 2.4274 & 2.1034 & 1.4663 \\
    1.5 & 2.4737 & 2.4237 & 2.2463 & 1.8933 & 1.3592 \\
    2.0 & 2.3237 & 2.2602 & 2.0639 & 1.7283 & 1.2586 \\
    2.5 & 2.1459 & 2.0830 & 1.8950 & 1.5862 & 1.1598 \\
    3.0 & 1.9721 & 1.9137 & 1.7404 & 1.4576 & 1.0655 \\
    \bottomrule
  \end{tabular}
\end{table}

由结果可见三点。

第一，\textbf{温度上升比问题一快得多}。问题一（附录 2）在 $1800$\,s 时中心仅
$33.58\,^\circ$C，而问题二已达 $40.34\,^\circ$C。原因是附录 3 的热容
$\rho c_p$ 在初期虽更大，但水分快速脱除使 $\rho c_p$ 显著下降
（$C$ 由 $2.55$ 降到 $0.15$ 时 $\rho c_p$ 由 $3.33\times10^6$ 降到
$1.21\times10^6$\,J/(m$^3\cdot$K)），热扩散率随之提高；同时式
\eqref{eq:p2-D} 中的 $e^{-3850/T}$ 使 $D$ 随升温急剧增大，
两个效应互相强化。

第二，\textbf{2.5\,h 后药材基本达到烘房温度}。$3$\,h 时中心 $49.83\,^\circ$C、
表面 $49.96\,^\circ$C，径向温差仅 $0.12\,^\circ$C。预热平衡阶段至此结束，
过程进入恒温干燥阶段。

第三，\textbf{水分脱除呈现出明显的"由表及里"次序}。$0.5$\,h 时中心仍是
$2.5500$，表面已降到 $1.6233$；$1$\,h 时中心才开始下降（$2.5435$）。
到 $3$\,h，表面降到 $1.0655$，而中心仍有 $1.9721$，径向梯度很大。
这说明过程始终由内部扩散控制。

\subsection{模型检验}
\label{sec:p23-verify}

\subsubsection{求解器交叉验证}

把附录 3 的物性替换为附录 2 的常数物性后，本模型应当退化为问题一的模型。
实测：$t=1800$\,s 时温度场与独立的 \texttt{solve\_p1.py} 结果最大偏差为
$2.6\times10^{-12}\,^\circ$C，达到机器精度，说明离散格式实现正确。

\subsubsection{网格与时间步收敛性}

表 \ref{tab:p23-conv} 给出问题三烘干时间随 $N$ 与 $\Delta t$ 的变化。
网格加密时结果单调收敛，$N=100$ 与 $N=200$ 相差 $0.16\%$；
时间步由 $5$\,s 减半到 $2$\,s 时烘干时间仅变化 $12$\,s（$0.005\%$），
说明本文采用的 $N=200$、$\Delta t=5$\,s 已充分收敛。

\begin{table}[htbp]
  \centering
  \caption{问题三烘干时间的收敛性}
  \label{tab:p23-conv}
  \small
  \begin{tabular}{ccccc}
    \toprule
    网格 / 时间步 & $N=50$ & $N=100$ & $N=200$ & $N=100,\ \Delta t=2$\,s \\
    \midrule
    烘干时间/s & 218575 & 219340 & 219685 & 219328 \\
    \bottomrule
  \end{tabular}
\end{table}

\section{烘干时间的确定（问题三）}

\subsection{问题分析}

烘干要求是"药材各处的\textbf{水分浓度应低于} $0.15$\,kg/kg"。
由于水分由表面向内部逐步脱除，整个过程中心处的含水率始终最高，
因此判据等价于
\begin{equation}
  t_{\mathrm{dry}}=\min\left\{t\ \middle|\
  \max_{0\le r\le R}C(r,t)<0.15\right\}
  =\min\left\{t\ \middle|\ C(0,t)<0.15\right\}.
  \label{eq:p3-criterion}
\end{equation}
数值上取求解过程中首次满足 $C(0,t)<0.15$ 的时刻。

\subsection{结果}

模型给出的烘干时长为
\begin{equation}
  t_{\mathrm{dry}}=219685\ \mathrm{s}=61.02\ \mathrm{h}=2.543\ \text{天}.
\end{equation}

该结果与题面"烘干过程一般持续 $2\sim3$ 天"的描述一致，
是对模型合理性的一个独立佐证。表 \ref{tab:p3} 给出每隔 $6$\,h 的
水分浓度分布。

\begin{table}[htbp]
  \centering
  \caption{药材烘干过程的水分浓度（单位：kg/kg）}
  \label{tab:p3}
  \begin{tabular}{cccccc}
    \toprule
    \multirow{2}{*}{时间/h} & \multicolumn{5}{c}{到药材中心的距离/cm} \\
    \cmidrule(lr){2-6}
     & 0     & 0.5   & 1.0   & 1.5   & 2.0 \\
    \midrule
    6  & 1.2031 & 1.1676 & 1.0616 & 0.8841 & 0.6206 \\
    12 & 0.5372 & 0.5215 & 0.4732 & 0.3864 & 0.2113 \\
    18 & 0.3328 & 0.3243 & 0.2976 & 0.2475 & 0.0977 \\
    24 & 0.2559 & 0.2502 & 0.2322 & 0.1976 & 0.0712 \\
    30 & 0.2174 & 0.2130 & 0.1992 & 0.1719 & 0.0621 \\
    36 & 0.1941 & 0.1905 & 0.1790 & 0.1561 & 0.0579 \\
    42 & 0.1784 & 0.1752 & 0.1652 & 0.1451 & 0.0556 \\
    48 & 0.1669 & 0.1641 & 0.1552 & 0.1370 & 0.0543 \\
    54 & 0.1581 & 0.1555 & 0.1474 & 0.1307 & 0.0534 \\
    60 & 0.1511 & 0.1487 & 0.1411 & 0.1257 & 0.0527 \\
    烘干结束 & 0.1500 & 0.1477 & 0.1402 & 0.1249 & 0.0527 \\
    \bottomrule
  \end{tabular}
\end{table}

\subsection{结果分析}

从表 \ref{tab:p3} 可以看出烘干过程的三个典型阶段。

\textbf{快速干燥段（$0\sim12$\,h）}：含水率大幅下降，中心由 $2.55$ 降到
$0.5372$。此阶段表面水分充足，内部扩散是限制环节，干燥速率基本恒定。

\textbf{降速干燥段（$12\sim36$\,h）}：速率明显放缓，中心由 $0.5372$ 降到
$0.1941$。表面含水率已接近空气平衡值（$0.05$ 附近），驱动力大幅衰减；
同时式 \eqref{eq:p2-D} 中 $e^{-0.45/C}$ 使 $D$ 随含水率降低而急剧减小，
内部水分越来越难迁移。

\textbf{极慢干燥段（$36\sim61$\,h）}：中心含水率由 $0.1941$ 缓缓逼近 $0.15$，
耗时约 $25$\,h，占全程的 $41\%$。这一段在工程上是最"不划算"的：
大量时间只换来很小的含水率下降。若对产品均匀性要求可以放宽，
适当延长或改用变温工艺会显著节能。

需要指出，判据由\textbf{中心}决定，而中心恰是最难干燥的位置。
$61.02$\,h 时有 $C(0)=0.1500$、$C(0.5\,\mathrm{cm})=0.1477$、
$C(2\,\mathrm{cm})=0.0527$，表面已干燥到接近平衡而中心刚刚达标，
内外含水率相差近三倍。

\section{考虑尺寸变化的烘干模型（问题四）}

\subsection{问题分析}

实际烘干中水分流失会引起药材收缩。附件 2 给出了半径随时间的实测变化：
$R$ 由 $2.000$\,cm 单调收缩至 $1.198$\,cm，收缩比 $0.599$，
对应截面积减为 $0.359$ 倍。

收缩带来了三个变化：其一，扩散路径变短，有利于干燥；
其二，表面几何面积减小，但对单位体积而言换热传质面积反而增大；
其三，计算区域随时间缩小，属于\textbf{动边界问题}，需要专门的坐标处理。

附件 2 的一个关键特征是收缩高度前置：
完成 $50\%$ 的收缩只需 $2.5$\,h，$90\%$ 只需 $10$\,h，$99\%$ 在 $22$\,h 内完成，
之后到 $72$\,h 半径几乎不再变化。也就是说，\textbf{收缩集中在干燥前期}，
而含水率的缓慢下降一直持续到 $50$\,h 之后。因此半径与含水率之间
不存在简单的线性关系，不宜用 $R=R_0 C/C_0$ 一类的经验式，
本文直接采用附件 2 的实测 $R(t)$。

\subsection{模型假设}

在问题二假设的基础上修改：

\begin{enumerate}
  \item[(C1)] 药材各点按比例均匀收缩，半径由附件 2 的 $R(t)$ 给定，
        即任意时刻的物理半径为 $r=\xi R(t)$，$\xi\in[0,1]$ 为贴体坐标；
  \item[(C2)] 物性改用附录 4 的经验公式
        \begin{equation}
          \rho=760+90C,\quad
          c_p=1850+2150\frac{C}{C+1},\quad
          k=0.12+0.20\frac{C}{C+1},
          \label{eq:p4-props}
        \end{equation}
        \begin{equation}
          D=4.2\times10^{-4}e^{-0.30/C}e^{-3850/T};
          \label{eq:p4-D}
        \end{equation}
  \item[(C3)] 初始半径 $R(0)=2.000$\,cm，与问题二一致，
        故初始条件仍为 $T=28\,^\circ$C、$C=2.55$\,kg/kg。
\end{enumerate}

\subsection{模型建立}

\subsubsection{贴体坐标变换}

令
\begin{equation}
  \xi=\frac{r}{R(t)}\in[0,1],
\end{equation}
把移动区域 $[0,R(t)]$ 映射为固定区域 $[0,1]$。由复合求导，
对固定物理位置 $r$ 求时间导数时
\begin{equation}
  \left.\frac{\partial}{\partial t}\right|_{r}
  =\left.\frac{\partial}{\partial t}\right|_{\xi}
   -\frac{\xi\dot R}{R}\frac{\partial}{\partial\xi},
  \qquad
  \frac{\partial}{\partial r}=\frac{1}{R}\frac{\partial}{\partial\xi},
  \label{eq:p4-transform}
\end{equation}
其中 $\dot R=\mathrm{d}R/\mathrm{d}t<0$。代入式 \eqref{eq:p2-heat}
并整理，得
\begin{equation}
  \rho c_p\,\xi R^2\frac{\partial T}{\partial t}
  =\frac{\partial}{\partial\xi}
   \left(k\,\xi\frac{\partial T}{\partial\xi}\right)
   +\rho c_p\,\xi R\dot R\,\frac{\partial T}{\partial\xi},
  \label{eq:p4-heat}
\end{equation}
水分方程同理
\begin{equation}
  \left(\rho+C\frac{\mathrm{d}\rho}{\mathrm{d}C}\right)\xi R^2
   \frac{\partial C}{\partial t}
  =\frac{\partial}{\partial\xi}
   \left(\rho D\,\xi\frac{\partial C}{\partial\xi}\right)
   +\left(\rho+C\frac{\mathrm{d}\rho}{\mathrm{d}C}\right)
    \xi R\dot R\,\frac{\partial C}{\partial\xi}.
  \label{eq:p4-mass}
\end{equation}
右端最后一项是坐标变换带来的\textbf{对流项}，物理上对应材料被
压缩（$\dot R<0$）对水分与热量的输运。若令 $\dot R\equiv0$，
式 \eqref{eq:p4-heat}、\eqref{eq:p4-mass} 退化为问题二的方程组。

\subsubsection{定解条件}

初始条件与中心对称条件同式 \eqref{eq:p2-ic}、\eqref{eq:p2-center}。
表面 $\xi=1$ 处边界条件形式不变：
\begin{equation}
  -k\left.\frac{\partial T}{\partial\xi}\right|_{\xi=1}
  =hR\left(T_s-T_\infty\right),\qquad
  -D\left.\frac{\partial C}{\partial\xi}\right|_{\xi=1}
  =h_mR\left(C_s-C_\infty\right).
  \label{eq:p4-bc}
\end{equation}

\subsubsection{$R(t)$ 的处理}

附件 2 为每 $1800$\,s 一个点的离散数据。若直接线性插值，
$\dot R$ 在节点处跳变，会把间断引入对流项。因此本文采用
\textbf{单调三次 Hermite 插值（PCHIP，Fritsch--Carlson 构造）}，
它既保证 $C^1$ 连续，又保持数据的单调性、不产生过冲。
收缩终止的时刻由数据本身决定，无需外推。

\subsection{数值实现}

空间离散仍用节点中心有限体积，但控制体权重变为
$w_i=\int\xi\,\mathrm{d}\xi$，即
\begin{equation}
  w_0=\frac{h^2}{8},\qquad w_i=\xi_i h\ (1\le i\le N-1),\qquad
  w_N=\frac{h(1-h/4)}{2}.
\end{equation}
对流项在内部节点取中心差分、在表面节点取单侧差分。
由于 $\dot R$ 在初期可达 $-2.9\times10^{-5}$\,m/s，对流项与扩散项量级相当，
不忽略等号右端第二项是必要的。时间方向与非线性处理同问题二，
取 $N=200$、$\Delta t=5$\,s。

\subsection{结果}

模型给出的烘干时长为
\begin{equation}
  t_{\mathrm{dry}}=181115\ \mathrm{s}=50.31\ \mathrm{h}=2.096\ \text{天},
\end{equation}
比不考虑收缩的问题三（$61.02$\,h）缩短了约 $17.6\%$。

\begin{table}[htbp]
  \centering
  \caption{药材烘干过程的水分浓度（含收缩，单位：kg/kg）}
  \label{tab:p4}
  \small
  \begin{tabular}{ccccccc}
    \toprule
    \multirow{2}{*}{时间/h} & \multicolumn{5}{c}{到药材中心的距离/cm}
      & \multirow{2}{*}{药材表面} \\
    \cmidrule(lr){2-6}
     & 0     & 0.5   & 1.0   & 1.5   & 2.0 & \\
    \midrule
    6  & 1.4770 & 1.2770 & 0.7859 & —— & —— & 0.2765 \\
    12 & 0.6278 & 0.5552 & 0.3453 & —— & —— & 0.1293 \\
    18 & 0.3709 & 0.3351 & 0.2201 & —— & —— & 0.0801 \\
    24 & 0.2694 & 0.2470 & 0.1707 & —— & —— & 0.0646 \\
    30 & 0.2185 & 0.2024 & 0.1450 & —— & —— & 0.0585 \\
    36 & 0.1884 & 0.1757 & 0.1292 & —— & —— & 0.0555 \\
    42 & 0.1686 & 0.1580 & 0.1185 & —— & —— & 0.0539 \\
    48 & 0.1544 & 0.1453 & 0.1106 & —— & —— & 0.0529 \\
    烘干结束 & 0.1500 & 0.1413 & 0.1081 & —— & —— & 0.0526 \\
    \bottomrule
  \end{tabular}
\end{table}

表 \ref{tab:p4} 中 $1.5$\,cm 与 $2.0$\,cm 两列以"——"表示，
原因是半径收缩至 $1.5$\,cm 以下之后（约 $5.5$\,h）这些位置已不在药材内部，
水分浓度无定义。这也是题面表 6 用"药材表面"而非固定距离作为末列的原因：
表面位置本身是随时间移动的。

\subsection{收缩对烘干时长的影响}

为分离收缩的贡献，用附录 4 的同一组物性但固定半径 $R\equiv2$\,cm 再求解一次，
得烘干时长 $493555$\,s $=5.71$ 天；计入收缩后为 $2.10$ 天，
缩短为原来的 $36.7\%$。

这一效应可从三个层面理解：扩散路程由 $R$ 缩短至 $0.6R$，
特征扩散时间按 $R^2$ 计减少到 $36\%$；
单位体积的表面积按 $1/R$ 计增大约 $1.67$ 倍，
表面交换能力相对增强；此外对流项 $\propto\dot R$ 把表层水分
向内部"挤压"，进一步加速了内部水分的再分布。

这一结论有明确的工程含义：\textbf{忽视收缩会严重高估烘干时长}
（本例中高估约 $1.7$ 倍）。若按不考虑收缩的结果安排生产，
会造成过度干燥、能耗浪费，甚至因干裂而降低药材品质。

\subsection{模型检验}

\begin{enumerate}
  \item \textbf{退化检验}：令 $\dot R\equiv0$，式 \eqref{eq:p4-heat}、
        \eqref{eq:p4-mass} 退化为问题二模型，程序实测退化解与问题二解
        完全一致，说明贴体坐标与对流项的实现正确。
  \item \textbf{单调性检验}：PCHIP 插值后的 $R(t)$ 单调递减，
        $\dot R\le0$ 始终成立，未出现插值过冲引起的非物理膨胀。
  \item \textbf{合理性检验}：收缩使烘干时长缩短，
        与"扩散路径变短有利于干燥"的物理预期一致；
        问题三 $61.02$\,h、问题四 $50.31$\,h 均落在题面所述 $2\sim3$ 天区间内。
\end{enumerate}

\section{结果文件说明}

四个结果文件均按附件 3 模板的格式保存，所有数值保留四位小数：

\begin{table}[htbp]
  \centering
  \caption{结果文件一览}
  \label{tab:files}
  \small
  \begin{tabular}{llll}
    \toprule
    文件 & 工作表 & 时间范围与间隔 & 规模 \\
    \midrule
    \texttt{result1.xlsx} & 温度、水分浓度 & $1\sim1800$\,s，每 $1$\,s & $1800\times21$ \\
    \texttt{result2.xlsx} & 温度、水分浓度 & $1\sim219685$\,s，每 $1$\,s & $219685\times21$ \\
    \texttt{result3.xlsx} & 水分浓度 & $60\sim219660$\,s，每 $60$\,s & $3661\times21$ \\
    \texttt{result4.xlsx} & 水分浓度 & $60\sim181080$\,s，每 $60$\,s & $3018\times21$\\
    \bottomrule
  \end{tabular}
\end{table}

各表的 A 列均为时间（单位：s），第 1 行为到药材中心的距离（单位：cm），
径向间隔 $0.1$\,cm。\texttt{result4.xlsx} 在二十一个固定位置之后
另加"药材表面"一列；因药材收缩，固定位置中超出当前半径者留空。

\end{document}
